% --- subsection1_2.tex ---
\subsection{Phân cụm danh từ (Noun Phrase Chunking)}

\subsubsection{Mục tiêu bài toán}
Sau khi các mệnh đề đã được tách rời ở Tác vụ 1.1, bước tiếp theo là xác định các thực thể hoặc khái niệm quan trọng thông qua việc phân cụm danh từ. Mục tiêu là nhóm các từ đơn lẻ có liên quan chặt chẽ về mặt ngữ pháp để tạo thành các cụm danh từ (Noun Phrases) hoàn chỉnh, giúp máy tính hiểu được các đối tượng cụ thể (ví dụ: "Hệ thống Quản lý Bệnh viện", "Bộ Y tế") thay vì các từ rời rạc.

\subsubsection{Phương pháp tiếp cận và Triển khai}
Nhóm tiếp tục sử dụng thư viện \texttt{py\_vncorenlp} để thực hiện gán nhãn từ loại (POS Tagging) và xây dựng thuật toán trích xuất dựa trên tập luật về từ loại. Quy trình gồm các bước sau:

\begin{enumerate}
    \item \textbf{Xác định tập từ loại mục tiêu (\texttt{is\_compound\_noun\_part}):} 
    Dựa trên đặc điểm ngôn ngữ học tiếng Việt, một cụm danh từ thường được cấu thành từ các từ loại sau:
    \begin{itemize}
        \item \texttt{N}: Danh từ chung.
        \item \texttt{Np}: Danh từ riêng.
        \item \texttt{Nc}: Danh từ chỉ loại.
        \item \texttt{Nu}: Danh từ đơn vị.
        \item \texttt{M}: Số từ (thường đi kèm danh từ để chỉ số lượng).
    \end{itemize}

    \item \textbf{Thuật toán trích xuất (\texttt{extract\_noun\_phrases}):} 
    Nhóm triển khai thuật toán duyệt tuyến tính qua danh sách các token đã được gán nhãn:
    \begin{itemize}
        \item Hệ thống sẽ gom nhóm các token liên tiếp nhau nếu chúng đều thuộc tập từ loại mục tiêu đã xác định ở trên.
        \begin{quote}
            \textit{Ví dụ:} "Hệ/N thống/N Quản/V lý/N" $\rightarrow$ "Hệ\_thống" (Do "Quản" là động từ nên sẽ ngắt cụm tại đó).
        \end{quote}
        \item Các từ trong cùng một cụm được kết nối với nhau bằng dấu gạch dưới (\texttt{\_}) để tạo thành một đơn vị ngữ nghĩa duy nhất.
        \item Chỉ các cụm có độ dài từ 2 token trở lên mới được coi là cụm danh từ phức hợp để trích xuất.
    \end{itemize}

    \item \textbf{Xử lý dữ liệu:}
    Dữ liệu đầu vào là file \texttt{clauses.txt} (kết quả từ Tác vụ 1.1). Thuật toán sẽ xử lý từng mệnh đề, trích xuất tất cả các cụm danh từ xuất hiện và lưu trữ chúng.
\end{enumerate}

\subsubsection{Kết quả đầu ra}
Kết quả cuối cùng được lưu tại file \texttt{noun\_phrases.txt}. Mỗi dòng trong file tương ứng với các cụm danh từ tìm thấy trong một mệnh đề, được phân cách rõ ràng. Việc phân cụm này tạo tiền đề quan trọng cho tác vụ trích xuất thực thể (NER) và phân tích quan hệ ở các giai đoạn sau của dự án.